\section{Related Work}
\label{sec:related}

\textbf{Parameter-Efficient Fine-Tuning and Model Merging.}
Parameter-Efficient Fine-Tuning (PEFT) techniques, such as LoRA \cite{hu2021lora}, Compacter \cite{karimi2022compacter}, and VL-Adapter \cite{sung2022vladapter}, have revolutionized multi-task adaptation by training small specialized parameter modules instead of full networks. To consolidate multiple expert capabilities without the serving overhead of running separate models, weight-space model merging has emerged as a key solution. Initial paradigms, such as Model Soups \cite{wortsman2022model} and Fisher-weighted averaging \cite{matena2021merging}, rely on simple linear combinations of expert weights. More advanced methods, such as Ties-Merging \cite{yadav2023ties}, address parameter interference, while AdaMerging \cite{adamerging} learns adaptive merging coefficients on calibration subsets. However, these paradigms are strictly \emph{static} and \emph{offline}; they compute a single set of merged weights once before deployment, leaving the model incapable of responding to test-time shifts.

\textbf{Test-Time Adaptation and Dynamic Ensembling.}
Test-time adaptation (TTA) seeks to adjust model parameters or activations on-the-fly to handle domain shifts or out-of-distribution (OOD) test inputs. Tent \cite{wang2020tent} minimizes entropy at test-time, while source-free domain adaptation methods \cite{liang2023source} leverage pseudo-labeling. In the context of multi-expert serving, dynamic model ensembling blends expert weights or activations dynamically for each query. For instance, SABLE \cite{wang2020tent} extracts coordinate similarity projections to route inputs locally to task centroids in a stateless fashion. While stateless dynamic ensembling provides high local plasticity, it is extremely sensitive to input noise, leading to severe high-frequency oscillations in routing decisions across adjacent layers.

\textbf{Stateful Serving and Temporal Routing.}
To mitigate the jitter of stateless routing on continuous task streams, recent literature has explored stateful routing paradigms that enforce temporal coherence. Momentum-Merge models sequential dependencies by applying a standard Exponential Moving Average (EMA) in Euclidean space over expert activations. More complex formulations, such as ChemMerge, draw inspiration from biochemical kinetics, formulating state updates as continuous-time ODEs where activation signals act as chemical concentration potentials. Although stateful, these approaches operate in flat Euclidean spaces $\mathbb{R}^K$ and project the resulting coordinates onto the probability simplex $\Delta^{K-1}$ post-hoc via Softmax or boundary clipping. This unconstrained-to-constrained mismatch results in representational hysteresis under task transitions, scale distortions of the feature activations, and requires costly virtual-time numerical ODE solvers.

\textbf{Non-Euclidean and Information-Geometric Deep Learning.}
Our work bridges dynamic model ensembling with non-Euclidean information geometry and Riemannian manifolds, while sharing deep mathematical connections with quantum probability. Operating on curved manifolds, such as the Grassmannian, Stiefel, or hyperspherical manifolds, has a rich history in robust representation learning \cite{gong2014multi}. Information-geometric and quantum-inspired neural networks utilize state representations on unit spheres or unit Hilbert spaces to capture rich, non-linear dependencies. Spherical linear interpolation (Slerp) was originally introduced by Shoemake \cite{shoemake1985animating} for smooth rotation animations and has since been adapted for generative model interpolations. However, UGR is the first to employ closed-form Fisher-Rao geodesic flow on the unit sphere $\mathbb{S}^{K-1}$ (corresponding to the information-geometric square-root mapping) to achieve exact, stable, and lag-free test-time model ensembling. By replacing unconstrained flat-space updates with curved manifold rotations, we establish a new geometric foundation for adaptive serving.
